\scanpage{8}
\diagram{\begin{tikzcd}[ampersand replacement=\&]
H^1(k,I) \arrow[r,dash] \arrow[d] \& H^1(k,I/I\cap K) \arrow[d] \\
H^1(k,\unsure{G|C}) \arrow[r,dash,"\unsure{f|j}_1"] \arrow[d,"u_1"'] \& H^1(k,G/K) \arrow[d,"v_1"] \\
H^1(k,G/I) \arrow[r,dash,"f'_1"] \& H^1(k,G/KI)
\end{tikzcd}}
Par l'hypothèse \inserted{a)} $u_1$ et $f'_1$ sont surjectifs, donc $f'_1u_1=v_1f_1$ l'est, donc si $\xi'\in H^1(k,G/K)$, il existe $\xi\in H^1(k,G)$ tel que $\xi'$ et $f_1(\xi)$ aient même image dans $H^1(k,G/KI)$. Quitte à tordre par $\xi'$, on peut supposer que $\xi'=0$, donc que $v_1(\xi)=0$, donc que $\xi$ est \deleted{dans} l'image d'un $\unsure{\chi|\zeta}\in H^1(k,KI/K)=H^1(k,I/I\cap K)$.

Or par hypothèse b), \deleted{\unsure{il s'ensuit}} $\chi$ est dans l'image de $H^1(k,I)$, et on gagne.

\uline{Lemme 4} Supposons qu'on ait \deleted{\unsure{a)}}

\unsure{c)|e)} $H^2(k,\unsure{\mathscr{K}|\mathcal{K}})=0$ si $K$ groupe \inserted{(commutatif)} fini radiciel sur $k$.

Alors pour tte \unsure{suite|schémas} exacte de groupes algébriques
\diagram{\begin{tikzcd}[ampersand replacement=\&]
e \arrow[r,dash] \& K \arrow[r,dash] \& G \arrow[r,dash] \& G' \arrow[r,dash] \& e
\end{tikzcd}}
avec $K$ \uline{radiciel} [pas néc. comm!] , on a
\[
H^1(k,G)-H^1(k,G')\quad\text{\uline{surjectif}}.
\]

Par dévissage, on peut supposer $K$ ``de hauteur 1''. Soit $\unsure{\mathfrak{k}|k}$ son Algèbre de
% archive sheet number: 7

